% docs/planteamiento.tex
\section{Planteamiento del problema}

En el departamento del Meta (Colombia), las condiciones ambientales agresivas —altas temperaturas, humedad elevada y particularidades de los suelos— inciden de manera directa en el desempeño y la durabilidad del concreto empleado en obras civiles. A pesar de ello, existe una carencia de herramientas digitales locales que apoyen, de forma didáctica y con trazabilidad, la dosificación y la selección de aditivos según el contexto regional. Esta brecha se hace evidente en los laboratorios de materiales y en los semilleros de investigación de la Universidad Cooperativa de Colombia (UCC), donde la replicabilidad de los ensayos y la sistematización de resultados son esenciales para la formación y la transferencia de conocimiento.

La ausencia de software especializado para apoyar el diseño y la verificación de mezclas adaptadas a condiciones locales puede derivar en decisiones inconsistentes, menor durabilidad y mayores costos de mantenimiento. De acuerdo con lineamientos nacionales, es prioritario impulsar la innovación tecnológica en el sector construcción para elevar la calidad y sostenibilidad de la infraestructura (Departamento Nacional de Planeación, 2018). A su vez, gremios y estudios académicos resaltan la necesidad de incorporar herramientas digitales que integren particularidades ambientales y estándares técnicos en el diseño de mezclas para optimizar recursos y asegurar la vida útil de las estructuras (CCI, 2020; González et al., 2019).

En este contexto, se identifica la necesidad de una solución \textbf{dirigida a los estudiantes del semillero de investigación y a quienes participan en las prácticas de laboratorio de materiales de la UCC}, que les permita: (i) registrar de manera estructurada las condiciones de diseño, (ii) asistir la dosificación base y la selección de aditivos con proyecciones de resistencia a distintas edades, y (iii) garantizar la trazabilidad de los ensayos y ejercicios académicos. Es importante precisar que \textbf{DIAPCE no está concebida como una herramienta de uso general ni para profesionales en ejercicio fuera del entorno universitario}; su alcance es intencionalmente formativo e institucional. DIAPCE responde a esta necesidad mediante una aplicación multiplataforma desarrollada en Flutter que opera sin conexión (SQLite), orientada exclusivamente a estudiantes del semillero y los laboratorios de la UCC para apoyar decisiones basadas en datos, estandarizar prácticas académicas y facilitar la comparación reproducible de resultados. La herramienta no sustituye el criterio del docente ni el marco normativo, sino que lo complementa con evidencia sistematizada y accesible en el contexto regional del Meta.
